\documentclass[12pt,a4paper]{article}
\usepackage[spanish,es-tabla]{babel}
\decimalpoint
\usepackage{array}
\newcolumntype{C}[1]{>{\centering\arraybackslash}m{#1}}
\usepackage{bbm}
\usepackage{tikz}
\usetikzlibrary{fit,positioning}
\usepackage{pgfplots}
\pgfplotsset{compat=1.18}
\usepackage{amsmath}
\usepackage{braket}
\usepackage[utf8]{inputenc}
\usepackage{graphicx}
\usepackage{subcaption}
\usepackage{cancel}
\usepackage{multicol}
\usepackage{wrapfig}
\usepackage[T1]{fontenc}
\usepackage{url}
\usepackage{graphicx}
\usepackage{gensymb}
\usepackage{booktabs}
\usepackage{amssymb, amsmath}
\usepackage{wrapfig}
\parskip=1 mm
\oddsidemargin -1.5 cm
\headsep -2 cm
\textwidth=17.8cm
\textheight=25.5cm
\begin{document}
\begin{center}
    {\LARGE \textbf{Introducción a la Mecánica Cuántica}} \\[0.4cm]

    {\large Facultad de Ciencias, UNAM} \\[0.6cm]

    \large \textbf{Tarea 3} \\[0.2cm]

    \begin{tabular}{l}
        \small Autor: Fernando Naed Ortega Montero\\
        \small Fecha: \today
    \end{tabular}
\end{center}



\section*{Preguntas Conceptuales.}

\begin{itemize}

\item ¿Son lo mismo $\bra{\hpi}$

\end{itemize}

\end{document}